\Opensolutionfile{ans}[ans/ansCD2D1-3]
\setcounter{section}{2}
\section{Giá trị lớn nhất, giá trị nhỏ nhất của hàm số}
\subsection{Kiến thức cần nắm}
\subsubsection{Định nghĩa}
	\begin{dn}
	Cho hàm số $y=f(x)$ xác định trên tập $K$. Khi đó:
$$M=\max\limits_K f(x)\Leftrightarrow\heva{&f(x)\leq M,\forall x\in K\\&\exists x_0\in K\colon f(x_0)=M.}
\quad m=\min\limits_K f(x)\Leftrightarrow\heva{&f(x)\geq m,\forall x\in K\\&\exists x_0\in K\colon f(x_0)=m.}$$
\end{dn}
\subsubsection{Định lí}
\begin{dl}
	Mọi hàm số liên tục trên một đoạn đều có giá trị lớn nhất, giá trị nhỏ nhất trên đoạn đó.
\end{dl}
\subsubsection{Các phương pháp tìm giá trị lớn nhất và nhỏ nhất của hàm số}
\textit{\textbf{Bài toán. Tìm GTLN GTNN của hàm số $\mathbf{y=f(x)}$ trên đoạn $\mathbf{K=[a;b]}$}.}\\
\textbf{Bước 1.} Hàm số đã cho xác định và liên tục trên đoạn $[a;b]$.
Tính $f'(x)$ và tìm những điểm $x_i$ sao cho tại đó có đạo hàm bằng $0$ hoặc liên tục nhưng không có đạo hàm.\\
\textbf{Bước 2.} Tính $f(a), f(b), f(x_i)$.\\
\textbf{Bước 3.} Kết luận: $\heva{&\max\limits_{[a;b]} f(x)=\max\left\{f(a); f(b); f(x_i)\right\}\\&\min\limits_{[a;b]} f(x)=\min\left\{f(a); f(b); f(x_i)\right\}}$.\\
\begin{note}
Nếu $K=(a;b)$ thì ở bước 2 ta không xét $f(a)$ và $f(b)$. Khi đó hàm số $f(x)$ có thể sẽ không có GTLN, GTNN trên $(a;b)$.\\
Nếu $y=f(x)$ đồng biến trên $[a;b]$ thì $\min\limits_{[a;b]} f(x)=f(a)$ và $\max\limits_{[a;b]} f(x)=f(b)$.\\
Nếu $y=f(x)$ nghịch biến trên $[a;b]$ thì $\min\limits_{[a;b]} f(x)=f(b)$ và $\max\limits_{[a;b]} f(x)=f(a)$.
\end{note}
\subsection{Phân loại và phương pháp giải bài tập}
\begin{dang}{Tìm GTLN – GTNN trên một đoạn, khoảng}\
\underline{Phương pháp:} Thực hiện theo các bước ở trên.
\end{dang}
\subsubsection{Các ví dụ}
\begin{vd}%Ví dụ 1:%[2D1B3-1] 
	Tim giá trị lớn nhất và giá trị nhỏ nhất của hàm số $y=x^3-3x^2-9x+2$ trên $[-2; 2]$.
	\loigiai{
		Ta có: $y'=3x^2-6x-9=0\Leftrightarrow\hoac{&x=-1\in(-2; 2)\\&x=3\notin(-2; 2).}$ \\
		$y(-2)=0; y(2)=-20; y(-1)=7$.\\
		Suy ra $\max\limits_{[-2; 2]} y=7$; $\min\limits_{[-2; 2]} y=-20$.}
\end{vd}
\begin{vd}%Ví dụ 2: %[2D1B3-1]
	Tìm giá trị lớn nhất và giá trị nhỏ nhất của hàm số $y=\dfrac{3x-1}{x-3}$ trên $[0;2]$.
	\loigiai{
		Hàm số đã cho xác định và liên tục trên đoạn $[0;2]$
		Ta có $y'=\dfrac{-8}{(x-3)^2}<0,\forall x\in[0;2]$.\\
		Tính $y(0)=\dfrac{1}{3}; y(2)=-5$.\\
		Suy ra $\max\limits_{[0;2]} y=\dfrac{1}{3}$ khi $x=0$, $\min\limits_{[0;2]} y=-5$ khi $x=2$.}
\end{vd}
\begin{vd}%Ví dụ 3: %[2D1B3-1]
	Tìm giá trị lớn nhất và giá trị nhỏ nhất của hàm số $y=\dfrac{x^2+3}{x-1}$ trên đoạn $[2;4]$.
	\loigiai{
		Hàm số $y=\dfrac{x^2+3}{x-1}$ liên tục trên đoạn $[2;4]$.\\
		Ta có $y'(x)=\dfrac{x^2-2x-3}{(x-1)^2}\Rightarrow y'=0\Leftrightarrow\hoac{&x=-1\notin(2; 4)\\&x=3\in(2; 4).}$ \\
		Tính $y'(2)=7; y'(4)=\dfrac{19}{3}; y'(3)=6$.\\
		Suy ra $\max\limits_{[2;4]} y=7$ khi $x=2$, $\min\limits_{[2;4]} y=6$ khi $x=3$.}
\end{vd}
\begin{vd}%Ví dụ 4: %[2D1K3-1]
	Gọi $M,m$ lần lượt là giá trị lớn nhất và nhỏ nhất của hàm số $y=x+\sqrt{4-x^2}$. Hãy tính $P=M+m$?
	\loigiai{
		Tập xác định: $\mathscr{D}=[-2;2]$. Ta có: $y'=1-\dfrac{x}{\sqrt{4-x^2}}=\dfrac{\sqrt{4-x^2}-x}{\sqrt{4-x^2}}$.\\
		$y'=0\Leftrightarrow\sqrt{4-x^2}=x\Leftrightarrow\heva{&x\geq 0\\&x=\pm\sqrt{2}}\Leftrightarrow x=\sqrt{2.}$ \\
		$y(\sqrt{2})=2\sqrt{2},y(2)=2,y(-2)=-2$.\\
		Vậy $M=2\sqrt{2},m=-2\Rightarrow P=2\sqrt{2}-2=2(\sqrt{2}-1)$.}
\end{vd}
\begin{vd}%Ví dụ 5: %[2D1B3-1]
	Tìm giá trị lớn nhất và giá trị nhỏ nhất của hàm số $y=(x-6)\sqrt{x^2+4}$ trên đoạn $[0;3]$.
	\loigiai{
		Hàm số đã cho xác định và liên tục trên đoạn $[0;3]$.\\
		Ta có: $y'=\dfrac{2x^2-6x+4}{\sqrt{x^2+4}},x\in[0;3]$.\\
		$y'=0\Leftrightarrow\hoac{&x=1\\&x=2.}$ \\
		Tính $y(1)=-5\sqrt{5}; y(0)=-12; y(2)=-8\sqrt{2}; y(3)=-3\sqrt{13}$.\\
		Suy ra $\max\limits_{[0;3]} y=-3\sqrt{13}$ khi $x=3$, $\min\limits_{[0;3]} y=-12$ khi $x=0$.\\
		Sử dụng Casio.\\
		Nhập MODE 7, $f(X)=(x-6)\sqrt{x^2+4}$.\\
		Start? $0$ End? $3$ Step? $\dfrac{1}{6}$. Kết luận.}
\end{vd}
\begin{vd}%Ví dụ 6: %[2D1B3-1]
	Tìm giá trị lớn nhất và giá trị nhỏ nhất của hàm số $y=\sqrt{-x^2+5x+6}$ trên đoạn $[-1;6]$.
	\loigiai{
		Hàm số đã cho xác định và liên tục trên đoạn $[-1;6]$.\\
		Ta có: $y'=\dfrac{-2x+5}{2\sqrt{-x^2+5x+6}}$; $y'=0\Leftrightarrow x=\dfrac{5}{2}\in[-1;6]$.\\
		$y(-1)=y(6)=0,y\left(\dfrac{5}{2}\right)=\dfrac{7}{2}$.\\
		Vậy $\max\limits_{[-1;6]} y=\dfrac{7}{2}$ khi $x=\dfrac{5}{2}$ và $\min\limits_{[-1;6]} y=0$ khi $x=-1,x=6$.}
\end{vd}
\begin{vd}%Ví dụ 7: %[2D1B3-1]
	Tìm giá trị lớn nhất và giá trị nhỏ nhất của hàm số $y=x-\sin 2x$ trên đoạn $\left[-\dfrac{\pi}{2};\pi\right]$.
	\loigiai{
		$y'=1-2\cos 2x=0\Leftrightarrow\cos 2x=\dfrac{1}{2}=\cos\dfrac{\pi}{3}\Leftrightarrow x=\pm\dfrac{\pi}{6}+k\pi$.\\
		Xét $x\in\left[\dfrac{-\pi}{2};\pi\right]$ ta được $x=\pm\dfrac{\pi}{6};x=\dfrac{5\pi}{6}$.\\
		$f\left(\dfrac{-\pi}{2}\right)=-\dfrac{\pi}{2}; f(\pi)=\pi; f\left(\dfrac{-\pi}{6}\right)=-\dfrac{\pi}{6}+\dfrac{\sqrt{3}}{2}; f\left(\dfrac{\pi}{6}\right)=\dfrac{\pi}{6}-\dfrac{\sqrt{3}}{2}; f\left(\dfrac{5\pi}{6}\right)=\dfrac{5\pi}{6}+\dfrac{\sqrt{3}}{2}$.\\
		Suy ra $\max\limits_{\left[-\dfrac{\pi}{2};\pi\right]} y=\dfrac{5\pi+3\sqrt{3}}{6}$; $\min\limits_{\left[-\dfrac{\pi}{2};\pi\right]} y=\dfrac{-\pi}{2}$.\\
		Sử dụng Casio.\\
		Chuyển chế độ sang Rad: Shift MODE 4.
		Nhập MODE 7, $f(X)=x-\sin 2x$.\\
		Start? $-\dfrac{\pi}{2}$ End? $\pi$ Step? $\dfrac{3\pi}{2}$. Kết luận.}
\end{vd}
\begin{vd}%Ví dụ 8: %[2D1B3-1]
	Tìm giá trị nhỏ nhất của hàm số $y=\left|x^3-9x^2+24x-68\right|$ trên đoạn $[-1;4]$.
	\loigiai{
		Bảng biến thiên của hàm số $y=x^3-9x^2+24x-68$ trên $[-1;4]$ 
		\begin{center}
			\begin{tikzpicture}
			\tkzTabInit[nocadre=false,lgt=1.2,espcl=2.5,deltacl=0.6]
			{$x$ /0.6,$y'$ /0.6,$y$ /2}
			{$-1$,$2$,$4$}
			\tkzTabLine{,+,$0$,-,}
			\tkzTabVar{-/$-102$, +/$-48$,-/$-52$}
			\end{tikzpicture}
		\end{center}
		Suy ra BBT của hàm số $y=\left|x^3-9x^2+24x-68\right|$ trên đoạn $[-1;4]$ là
		\begin{center}
			\begin{tikzpicture}
			\tkzTabInit[nocadre=false,lgt=1.2,espcl=2.5,deltacl=0.6]
			{$x$ /0.6,$y'$ /0.6,$y$ /2}
			{$-1$,$2$,$4$}
			\tkzTabLine{,-,$0$,+,}
			\tkzTabVar{+/$102$, -/$48$,+/$52$}
			\end{tikzpicture}
		\end{center}
		Vậy GTNN của hàm số $y=\left|x^3-9x^2+24x-68\right|$ trên đoạn $[-1;4]$ bằng $48$.\\
		Sử dụng Casio.\\
		Nhập MODE 7, $f(X)=\left|x^3-9x^2+24x-68\right|$.\\
		Start? $-1$ End? $4$ Step? $\dfrac{5}{18}$. Kết luận.}
\end{vd}
\begin{vd}%Ví dụ 10: %[2D1B3-2] 
	Tìm giá trị lớn nhất, nhỏ nhất của hàm số $y=x-5+\dfrac{1}{x}$ trên khoảng $(0;+\infty)$.
	\loigiai{
		Ta có: $y'=1-\dfrac{1}{x^2}=\dfrac{x^2-1}{x^2}$ trên khoảng $(0;+\infty)$.\\
		$y'=0\Leftrightarrow\hoac{&x=1\in(0;+\infty)\\&x=-1\notin(0;+\infty).}$ \\
		Dựa vào BBT: 
	   \begin{center}
	   	\begin{tikzpicture}
	   	\tikzset{double style/.append style={double distance=1.5pt}}
	   	\tkzTabInit[nocadre=false,lgt=1.2,espcl=2.5,deltacl=0.6]
	   	{$x$ /0.6,$y'$ /0.6,$y$ /2}
	   	{$-\infty$,$-1$,$0$,$1$,$+\infty$}
	   	\tkzTabLine{,+,$0$,-,d,-,$0$,+,}
	   	\tkzTabVar{-/$-\infty$,+/$-7$,-D+/$-\infty$/$+\infty$,-/$-3$,+/$+\infty$}
	   	\end{tikzpicture}
	   \end{center}
		Trên khoảng $(0;+\infty)$, không có $GTLN$; $GTNN=-3$.}
\end{vd}
\begin{vd}%Ví dụ 11: %[2D1B3-2]
	Tìm giá trị lớn nhất và giá trị nhỏ nhất của hàm số $y=x^2+\dfrac{2}{x}$ trên khoảng $(0;+\infty)$.
	\loigiai{
		Ta có: $\lim\limits_{x\to 0^+} y=+\infty$; $\lim\limits_{x\to+\infty} y=+\infty$.\\
		$y'=2x-\dfrac{2}{x^2}$.\\
		$y'=0\Leftrightarrow x=1$.\\
		Bảng biến thiên
		 \begin{center}
			\begin{tikzpicture}
			\tikzset{double style/.append style={double distance=1.5pt}}
			\tkzTabInit[nocadre=false,lgt=1.2,espcl=2.5,deltacl=0.6]
			{$x$ /0.6,$y'$ /0.6,$y$ /2}
			{$0$,$1$,$+\infty$}
			\tkzTabLine{d,-,$0$,+,}
			\tkzTabVar{D+/$-\infty$/$+\infty$,-/$-3$,+/$+\infty$}
			\end{tikzpicture}
		\end{center}
		Vậy không tồn tại giá trị lớn nhất, $\min\limits_{(0;+\infty)} y=3$.}
\end{vd}
\begin{vd}%Ví dụ 12: %[2D1K3-2]
	Tìm giá trị lớn nhất và giá trị nhỏ nhất của hàm số $y=\dfrac{3x^2+2x+3}{x^2+1}$.
	\loigiai{
		Tập xác định $\mathscr{D}=\mathbb{R}$.\\
		Ta có $y'=\dfrac{-2x^2+2}{\left(x^2+1\right)^2}$, $y'=0\Leftrightarrow-2x^2+2=0\Leftrightarrow\hoac{&x=1\\&x=-1.}$ \\
		$\lim\limits_{x\to\pm\infty} y=\lim\limits_{x\to\pm\infty}\dfrac{3x^2+2x+3}{x^2+1}=3\Rightarrow y=3$ là tiệm cận ngang.\\
		Bảng biến thiên: 
		\begin{center}
			\begin{tikzpicture}
			\tkzTabInit[nocadre=false,lgt=1.2,espcl=2.5,deltacl=0.6]
			{$x$ /0.6, $y'$ /0.6, $y$ /2.5}
			{$-\infty$,$-1$,$1$,$+\infty$}
			\tkzTabLine{,-,$0$,+,$0$,-,}
			\tkzTabVar{-/$3$,+/$2$,-/$4$,+/$3$}
			\end{tikzpicture}
		\end{center}
		Suy ra $\max\limits_{\mathbb{R}} y=4$ khi $x=1$, $\min\limits_{\mathbb{R}} y=2$ khi $x=-1$.}
\end{vd}
\subsubsection{Câu hỏi trắc nghiệm}
\begin{ex}%Câu 1.%[2D1B3-1]
	Giá trị nhỏ nhất $m$ của hàm số $y=x^3+6x^2-3$ trên đoạn $[-2;2]$ là
	\choice
	{$m=29$}
	{$m=13$}
	{\True $m=-3$}
	{$m=-4$}
\end{ex}
\begin{ex}%Câu 2.%[2D1B3-1]
	Giá trị nhỏ nhất của hàm số $y=\dfrac{1}{3}x^3-4x+2$ trên đoạn $[-1;3]$ là
	\choice
	{$y(-1)$}
	{\True $y(2)$}
	{$y(3)$}
	{$y(0)$}
\end{ex}
\begin{ex}%Câu 3.%[2D1B3-1]
	Tìm giá trị nhỏ nhất $M$ của hàm số $y=x^3-3x^2$ trên đoạn $[-1;1]$ 
	\choice
	{$M=-2$}
	{$M=0$}
	{\True $M=-4$}
	{$M=2$}
\end{ex}
\begin{ex}%Câu 4.%[2D1B3-1]
	Giá trị lớn nhất của hàm số $y=x^3-3x+3$ trên đoạn $\left[-3;\dfrac{3}{2}\right]$ là 
	\choice
	{$2$}
	{$3$}
	{$4$}
	{\True $5$}
\end{ex}
\begin{ex}%Câu 5.%[2D1B3-1]
	Giá trị lớn nhất của hàm số $y=x^3-2x^2-4x+5$   trên đoạn $[1;3]$ bằng
	\choice
	{$-3$}
	{$0$}
	{\True $2$}
	{$3$}
\end{ex}
\begin{ex}%Câu 6.%[2D1B3-1]
	Tìm giá trị nhỏ nhất $m$ của hàm số $y=x^4-x^2+13$ trên đoạn $[-2;3]$. 
	\choice
	{$m=\dfrac{51}{2}$}
	{$m=13$}
	{\True $m=\dfrac{51}{4}$}
	{$m=\dfrac{49}{4}$}
\end{ex}
\begin{ex}%Câu 7.%[2D1B3-1]
	Tìm giá trị lớn nhất $M$ của hàm số $y=x^4-2x^2+3$ trên đoạn $[0;\sqrt{3}]$ 
	\choice
	{$M=9$}
	{$M=8\sqrt{3}$}
	{\True $M=6$}
	{$M=1$}
\end{ex}
\begin{ex}%Câu 8.%[2D1B3-1]
	Tìm giá trị nhỏ nhất của hàm số $y=x^4-2x^2+3$ trên đoạn $[-3;2]$?
	\choice
	{$11$}
	{$0$}
	{$1$}
	{\True $2$}
\end{ex}
\begin{ex}%Câu 9.%[2D1B3-1]
	Tìm giá trị lớn nhất của hàm số $y=\dfrac{x-2}{x+1}$ trên đoạn $[0;2]$ 
	\choice
	{Không tồn tại}
	{\True $0$}
	{$-2$}
	{$2$}
\end{ex}
\begin{ex}%Câu 10.%[2D1B3-1]
	Tìm giá trị nhỏ nhất của hàm số $y=\dfrac{1-x}{2x-3}$ trên $[0;1]$. 
	\choice
	{$\min\limits_{[0;1]} y=0$}
	{\True $\min\limits_{[0;1]} y=-\dfrac{1}{3}$}
	{$\min\limits_{[0;1]} y=-1$}
	{$\min\limits_{[0;1]} y=-2$}
\end{ex}
\begin{ex}%Câu 11.%[2D1B3-1]
	Tìm giá trị lớn nhất của hàm số $y=\dfrac{3x-1}{x-3}$ trên đoạn $[0;2]$. 
	\choice
	{\True $\dfrac{1}{3}$}
	{$5$}
	{$-\dfrac{1}{3}$}
	{$-5$}
\end{ex}
\begin{ex}%Câu 12.%[2D1B3-1]
	Cho hàm số $y=\dfrac{x+1}{x-2}$. Giá trị nhỏ nhất của hàm số trên đoạn $[-4,0]$. 
	\choice
	{$y(-3)$}
	{\True $y(0)$}
	{$y(-1)$}
	{$y(-4)$}
\end{ex}

\begin{ex}%Câu 1.%[2D1B3-1]
	Hàm số $y=-x^3+3x-2$ trên đoạn $[-3; 0]$ có giá trị lớn nhất $M$, giá trị nhỏ nhất $m$. Tính giá trị của $M+m$. 
	\choice
	{$-6$}
	{$12$}
	{\True $14$}
	{$16$}
\end{ex}
\begin{ex}%Câu 2.%[2D1B3-1]
	Tính tổng giá trị lớn nhất và giá trị nhỏ nhất của hàm số $f(x)=x^3+3x^2-9x+1$ trên đoạn $[-4;0]$. 
	\choice
	{$24$}
	{\True $21$}
	{$22$}
	{$29$}
\end{ex}
\begin{ex}%Câu 3.%[2D1B3-1]
	Giá trị nhỏ nhất của hàm số $y=2x^3+3x^2-12x+2$ trên đoạn $[-1;2]$ đạt tại $x=x_0$. Giá trị $x_0$ bằng
	\choice
	{$1$}
	{$2$}
	{$-2$}
	{\True $-1$}
\end{ex}
\begin{ex}%Câu 4.%[2D1B3-1]
	Trên đoạn $[-2;2]$, hàm số nào sau đây có giá trị nhỏ nhất và giá trị lớn nhất trùng với giá trị cực đại và giá trị cực tiểu của nó?
	\choice
	{\True $y=x^3+2x-7$}
	{$y=x^3-6x-1$}
	{$y=x^3-2x^2+5x-10$}
	{$y=5x^3$}
\end{ex}
\begin{ex}%Câu 5.%[2D1B3-2]
	Tìm giá trị nhỏ nhất của hàm số $y=\left(x^2-2x+3\right)^2-7$. 
	\choice
	{Không tồn tại giá trị nhỏ nhất của hàm số}
	{$\min\limits_{\mathbb{R}} y=-5$}
	{$\min\limits_{\mathbb{R}} y=-7$}
	{\True $\min\limits_{\mathbb{R}} y=-3$}
\end{ex}
\begin{ex}%Câu 6.%[2D1B3-2]
	Giá trị nhỏ nhất của hàm số $y=x^3-3x+1$ trên khoảng $(0;2)$ là
	\choice
	{$3$}
	{$1$}
	{\True $-1$}
	{$0$}
\end{ex}
\begin{ex}%Câu 7.%[2D1B3-1]
	Tìm giá trị nhỏ nhất của hàm số $y=x^2+\dfrac{2}{x}$ trên đoạn $[2;3]$. 
	\choice
	{$\dfrac{15}{2}$}
	{$\dfrac{29}{3}$}
	{$3$}
	{\True $5$}
\end{ex}
\begin{ex}%Câu 8.%[2D1B3-1]
	Tìm giá trị nhỏ nhất của hàm số $y=x+\dfrac{4}{x}$ trên đoạn $[1; 4]$.
	\choice
	{\True $4$}
	{$2\sqrt{3}$}
	{$5$}
	{$3\sqrt{2}$}
\end{ex}
\begin{ex}%Câu 9.%[2D1B3-1]
	Tìm giá trị nhỏ nhất của hàm số $y=\dfrac{x^2+3}{x-1}$ trên $[2;4]$. 
	\choice
	{$\min\limits_{[2;4]} y=-2$}
	{\True $\min\limits_{[2;4]} y=6$}
	{$\min\limits_{[2;4]} y=-3$}
	{$\min\limits_{[2;4]} y=\dfrac{19}{3}$}
\end{ex}
\begin{ex}%Câu 10.%[2D1B3-1]
	Tìm giá trị nhỏ nhất của hàm số $y=x^2+\dfrac{2}{x}$ trên đoạn $[2;3]$. 
	\choice
	{$\dfrac{15}{2}$}
	{$\dfrac{29}{3}$}
	{$3$}
	{\True $5$}
\end{ex}
\begin{ex}%Câu 11.%[2D1B3-1]
	Tích của giá trị nhỏ nhất và giá trị lớn nhất của hàm số $f(x)=x+\dfrac{4}{x}$ trên $[1;3]$ bằng
	\choice
	{$\dfrac{52}{3}$}
	{\True $20$}
	{$6$}
	{$\dfrac{65}{3}$}
\end{ex}
\begin{ex}%Câu 12.%[2D1B3-1]
	Cho hàm số $f(x)=\dfrac{x^2-4x+7}{x-1}$. Gọi $M$, $m$ lần lượt là giá trị lớn nhất và giá trị nhỏ nhất của hàm số trên đoạn $[2; 4]$. Tính $M+m$?
	\choice
	{$M+m=7$}
	{$M+m=\dfrac{16}{3}$}
	{$M+m=\dfrac{13}{3}$}
	{\True $M+m=5$}
\end{ex}
\begin{ex}%Câu 13.%[2D1B3-1]
	Giá trị nhỏ nhất của hàm số $y=\dfrac{2x^2+5x+4}{x+2}$ trên đoạn $[0;1]$ là 
	\choice
	{\True $2$}
	{$3$}
	{$4$}
	{$5$}
\end{ex}
\begin{ex}%Câu 14.%[2D1B3-1]
	Giá trị lớn nhất của hàm số $y=\dfrac{2x^2+5x+8}{x+8}$ trên đoạn $[0;8]$ là 
	\choice
	{$12$}
	{\True $11$}
	{$10$}
	{$9$}
\end{ex}
\begin{ex}%Câu 15.%[2D1B3-1]
	Giá trị nhỏ nhất của hàm số $y=x+\dfrac{1}{x}$ trên khoảng $(0;+\infty)$ là 
	\choice
	{\True $2$}
	{$3$}
	{$4$}
	{$5$}
\end{ex}
\begin{ex}%Câu 16.%[2D1B3-1]
	GTNN của hàm số $f(x)=2\sin 2x-5x+1$ trên đoạn $\left[0;\dfrac{\pi}{2}\right]$ bằng 
	\choice
	{$3-\dfrac{5\pi}{4}$}
	{$0$}
	{$1$}
	{\True $1-\dfrac{5\pi}{2}$}
\end{ex}
\begin{ex}%Câu 17.%[2D1B3-1]
	Gọi $M$ và $m$ theo thứ tự là giá trị lớn nhất và nhỏ nhất của hàm số $y=x+\cos^2x$ trên đoạn $\left[0;\dfrac{\pi}{4}\right]$. Khi đó tổng $(M+m)$ bằng
	\choice
	{$-\dfrac{1}{2}$}
	{\True $\dfrac{\pi}{4}+\dfrac{3}{2}$}
	{$\dfrac{\pi}{4}+\dfrac{\pi}{6}$}
	{$\dfrac{\pi}{2}+\dfrac{3}{4}$}
\end{ex}
\begin{ex}%Câu 18.%[2D1K3-1]
	Tìm giá trị lớn nhất của hàm số $f(x)=\sqrt{x-2}+\sqrt{4-x}$. 
	\choice
	{$\sqrt{2}$}
	{\True $2$}
	{$4$}
	{$3$}
\end{ex}
\begin{ex}%Câu 19.%[2D1K3-1]
	Giá trị lớn nhất của hàm số $y=\sqrt{4-x^2}$ 
	\choice
	{\True $2$}
	{$3$}
	{$4$}
	{$5$}
\end{ex}
\begin{ex}%Câu 20.%[2D1K3-1]
	Gọi $M$, $m$ lần lượt là GTLN, GTNN của hàm số $y=\sqrt{x-1}+\sqrt{7-x}$. Khi đó có bao nhiêu số nguyên nằm giữa $M$, $m$?
	\choice
	{\True $1$}
	{$2$}
	{Vô số}
	{$0$}
\end{ex}
\begin{ex}%Câu 21.%[2D1K3-1]
	Cho hàm số $f(x)=\sqrt{2+x}+\sqrt{2-x}$. Khẳng định nào sau đây \textbf{đúng}?
	\choice
	{\True Giá trị lớn nhất của hàm số bằng $2\sqrt{2}$}
	{Hàm số đạt giá trị nhỏ nhất tại $x=0$}
	{Giá trị nhỏ nhất của hàm số bằng $0$}
	{Hàm số đạt giá trị lớn nhất tại $x=2$}
\end{ex}
\begin{ex}%Câu 22.%[2D1B3-1]
	Tìm giá trị lớn nhất của hàm số $y=x+\sqrt{2-x}$ trên đoạn $[-2;1]$ 
	\choice
	{$0$}
	{\True $2$}
	{$\dfrac{9}{4}$}
	{$\sqrt{2}$}
\end{ex}

\begin{ex}%Câu 1.%[2D1K3-1]
	Tìm giá trị lớn nhất của hàm số $f(x)=\dfrac{x+1}{\sqrt{x^2+1}}$?
	\choice
	{$1$}
	{$2$}
	{\True $\sqrt{2}$}
	{Không tồn tại}
	\loigiai{
		Tập xác định: $\mathscr{D}=\mathbb{R}$
		$f'(x)=\dfrac{\sqrt{x^2+1}-(x+1)\dfrac{2x}{2\sqrt{x^2+1}}}{x^2+1}=\dfrac{x^2+1-x^2-x}{\left(x^2+1\right)\sqrt{x^2+1}}=\dfrac{1-x}{\left(x^2+1\right)\sqrt{x^2+1}}$
		$f'(x)=0\Leftrightarrow x=1$, $f(1)=\sqrt{2}$. 
		\begin{center}
			\begin{tikzpicture}
			\tkzTabInit[nocadre=false,lgt=1.2,espcl=2.5,deltacl=0.6]
			{$x$ /0.6,$y'$ /0.6,$y$ /2}
			{$-\infty$,$1$,$+\infty$}
			\tkzTabLine{,+,$0$,-,}
			\tkzTabVar{-/$-\infty$, +/$\sqrt{2}$,-/$-\infty$}
			\end{tikzpicture}
		\end{center}
		Vậy giá trị lớn nhất của hàm số $f(x)=\dfrac{x+1}{\sqrt{x^2+1}}$ là $f(1)=\sqrt{2}$.}
\end{ex}
\begin{ex}%Câu 2.%[2D1K3-1]
	Giá trị lớn nhất của hàm số $y=\dfrac{x+\sqrt{1+9x^2}}{8x^2+1},x>0$. 
	\choice
	{$\dfrac{2\sqrt{2}}{3}$}
	{$\dfrac{\sqrt{2}}{3}$}
	{$\sqrt{2}$}
	{\True $\dfrac{3\sqrt{2}}{4}$}
	\loigiai{
		Hàm số đã cho xác định và liên tục trên khoảng $(0\colon +\infty)$
		$y=\dfrac{x+\sqrt{1+9x^2}}{8x^2+1}=\dfrac{9x^2+1-x^2}{\left(8x^2+1\right)\left(\sqrt{9x^2+1}-x\right)}=\dfrac{1}{\sqrt{9x^2+1}-x}$.\\
		Hàm số đạt giá trị lớn nhất trên khoảng $(0\colon +\infty)$ khi hàm số.\\
		$f(x)=\sqrt{9x^2+1}-x$ đạt giá trị nhỏ nhất trên khoảng $(0\colon +\infty)$.\\
		$f'(x)=\dfrac{9x}{\sqrt{9x^2+1}}-1$.\\
		$f'(x)=0\Leftrightarrow\sqrt{9x^2+1}=9x\Leftrightarrow\heva{&x>0\\&72x^2=1}\Leftrightarrow x=\dfrac{1}{6\sqrt{2}.}$ \\
		Vậy $\max\limits_{x>0} y=\dfrac{1}{\dfrac{2\sqrt{2}}{3}}=\dfrac{3\sqrt{2}}{4}$ khi $x=\dfrac{1}{6\sqrt{2}}$.}
\end{ex}
\begin{ex}%Câu 3.%[2D1K3-1]
	Tìm giá trị nhỏ nhất của hàm số $y=-x+3-\dfrac{1}{x+2}$ trên nửa khoảng $[-4;-2)$. 
	\choice
	{$\min \limits_{[-4;2)}y=4$}
	{$\min \limits_{[-4;2)}y=5$}
	{$\min \limits_{[-4;2)}y=\dfrac{15}{2}$}
	{\True $\min \limits_{[-4;2)}y=7$}
	\loigiai{
		TXĐ: $\mathscr{D}=[-4;-2)$; $y'=-1+\dfrac{1}{(x+2)^2}$; $y'=0\Rightarrow\hoac{&x=-1\notin[-4;-2)\\&x=-3\in[-4;-2).}$ \\
		Ta có
		\begin{center}
			\begin{tikzpicture}
			\tkzTabInit[nocadre=false,lgt=1.2,espcl=2.5,deltacl=0.6]
			{$x$ /0.6,$y'$ /0.6,$y$ /2}
			{$-4$,$3$,$2$}
			\tkzTabLine{,-,$0$,+,}
			\tkzTabVar{+/$\dfrac{15}{2}$, -/$7$,+/$+\infty$}
			\end{tikzpicture}
		\end{center}
		$y(-4)=\dfrac{15}{2};y(-3)=7;\lim\limits_{x\to-2^-} y=+\infty$.\\ Vậy $\min \limits_{[-4;2)}y=7$.}
\end{ex}
\begin{ex}%Câu 4.%[2D1K3-2]
	Giá trị nhỏ nhất của hàm số $y=x+2+\dfrac{1}{x-1}$ trên khoảng $(1;+\infty)$ là 
	\choice
	{$2$}
	{$3$}
	{$4$}
	{\True $5$}
	\loigiai{
		TXĐ: $\mathscr{D}=(1;+\infty)$; $y'=1-\dfrac{1}{(x-1)^2}$; $y'=0\Rightarrow\hoac{&x=0\notin[(1;+\infty)\\&x=2\in(1;+\infty).}$ \\
	Ta có
	\begin{center}
		\begin{tikzpicture}
		\tkzTabInit[nocadre=false,lgt=1.2,espcl=2.5,deltacl=0.6]
		{$x$ /0.6,$y'$ /0.6,$y$ /2}
		{$1$,$2$,$+\infty$}
		\tkzTabLine{d,-,$0$,+,}
		\tkzTabVar{D+/$+\infty$, -/$5$,+/$+\infty$}
		\end{tikzpicture}
	\end{center}
	Vậy $\min \limits_{(1;+\infty)}y=5$.}
\end{ex}
\begin{ex}%Câu 5.%[2D1K3-1]
	Tổng giá trị lớn nhất và nhỏ nhất của hàm số $y=f(x)=|x|+3$ trên $[-1; 1]$ là 
	\choice
	{$3$}
	{\True $7$}
	{$0$}
	{$4$}
	\loigiai{
		Ta có $y=f(x)=|x|+3=\hoac{&x+3 khi x\in[0; 1]\\&-x+3 khi x\in[-1; 0)}$ \\
		$ \Rightarrow\hoac{&y'=1 khi x\in[0; 1]\\&y'=-1 khi x\in[-1; 0).} $ \\
		Ta có bảng biến thiên của hàm số đã cho
		\begin{center}
			\begin{tikzpicture}
			\tkzTabInit[nocadre=false,lgt=1.2,espcl=2.5,deltacl=0.6]
			{$x$ /0.6,$y'$ /0.6,$y$ /2}
			{$-1$,$0$,$1$}
			\tkzTabLine{,-,d,+,}
			\tkzTabVar{+/$4$, -/$3$,+/$4$}
			\end{tikzpicture}
		\end{center}
		Vậy $\max \limits_{[-1; 1]}y+\min \limits_{[-1; 1]}y=7$.}
\end{ex}
\begin{ex}%Câu 6.%[2D1K3-1]
	Tìm câu sai trong các mệnh đề sau về GTLN và GTNN của hàm số $y=\left|x^3-3x+1\right|, x\in[0;3]$ 
	\choice
	{\True $\min y=1$}
	{$\max y=19$}
	{Hàm số có GTLN và GTNN}
	{Hàm số đạt GTLN khi $x=3$}
	\loigiai{
		$\max y=19$.}
\end{ex}
\begin{ex}%Câu 7.%[2D1K3-1]
	Gọi $M,N$ lần lượt là giá trị lớn nhất, giá trị nhỏ nhất của hàm số $f(x)=|x-3|\sqrt{x+1}$ trên đoạn $[0;4]$. Tính $M+2N$. 
	\choice
	{\True $\dfrac{16\sqrt{3}}{9}$}
	{$3+\sqrt{5}$}
	{$\dfrac{16\sqrt{3}}{3}$}
	{$\sqrt{5}$}
	\loigiai{		
		$f(x)=|x-3|\sqrt{x+1}=\sqrt{(x-3)^2(x+1)}$. Xét hàm số $g(x)=(x-3)^2(x+1),x\in[0;4]$.\\
		$g'(x)=2(x-3)(x+1)+(x-3)^2=(x-3)(2(x+1)+x-3)=(x-3)(3x-1)$.\\
		$g'(x)=0\Leftrightarrow\hoac{&x=3\in (0;4)\\&x=\dfrac{1}{3}\in (0;4)}$; $g(0)=9;\quad g\left(\dfrac{1}{3}\right)=\dfrac{256}{27};\quad g(3)=0;\quad g(4)=5$; $\max\limits_{[0;4]} g(x)=\dfrac{256}{27}$; $\min\limits_{[0;4]} g(x)=0$; $M=\dfrac{16\sqrt{3}}{9};N=0$.\\
		Vậy $M+2N=\dfrac{16\sqrt{3}}{9}$.}
\end{ex}
\Closesolutionfile{ans}